\documentclass[12pt]{article}

\usepackage{fontawesome}
\usepackage{hyperref}
\usepackage{xurl}
\usepackage[tagged, highstructure]{accessibility}
\usepackage{bookmark}
\usepackage{enumitem}

\hypersetup{
    colorlinks=false,
    pdfborder={0 0 0},
    pdftitle={Backup, Restore, and Post Mortem Lab},
    pdfauthor={Adrianna Holden-Gouveia},
    pdfsubject={Linux Administration - Lab Assignment},
    pdflang={en-US},
    pdfkeywords={backups, tar, cron, restore, incident response, post mortem}
}

% Configure PDF bookmarks for navigation
\bookmarksetup{
    numbered,
    open,
}

% Configure list spacing for better accessibility
\setlist{nosep}



\title{Backup, Restore, and Post Mortem}
\author{
        Adrianna Holden-Gouveia \\
        Website: \href{https://aholdengouveia.name}{aholdengouveia.name}\\ 
        \date{\vspace{-5ex}}
        %Email: \href{mailto:admin@aholdengouveia.name}{admin@aholdengouveia.name} \\
%        \faLinkedin{: aholdengouveia} \\
        \faGithub {: aholdengouveia} \\
       % \faTwitter {: aholdengouveia} \\
        }

%S\date{\today}


\begin{document}    

\maketitle

%\begin{abstract}
%This is a template for Linux Administration Lab work
%\end{abstract}
%\tableofcontents


\section*{Accessibility Notice}
This document is also available in HTML format at:

Link: \href{https://aholdengouveia.name/LinuxAdmin/labexcercises/backups.html}{\textbf{\nolinkurl{https://aholdengouveia.name/LinuxAdmin/labexcercises/backups.html}}}

The HTML version provides enhanced accessibility features including keyboard navigation, screen reader support, responsive design, dark mode support, and high contrast options.

\section*{Objectives:}
\begin{itemize}
    \item The objective of this lab is to give students practice with backups, disaster recovery, and incident response. Students will script and schedule backups, store them on a second server, prove they can restore from them, and write a post mortem. A backup is only useful if you know you can restore from it.
\end{itemize}

\section*{Background}
After last week's incident, management wants proof that Acme Corp can recover. A backup no one has tested isn't a backup. You will be working on the two servers you set up in the Server Setup lab: your Red Hat based server and your second server.

\section*{Requirements}
\begin{enumerate}
    \item Write a backup script using tar that backs up /etc, users' home directories, and your change log, compressed, with the date in the filename.
    \item Keep only the last 7 backups, deleting older ones automatically.
    \item Copy each backup to your other server, so it isn't stored only on the machine it came from.
    \item Schedule your backup script with cron.
    \item Restore drill: delete something important, such as a config file or a user's home directory, then restore it from your backup. Time how long it takes.
    \item Write a post mortem for either your restore drill or the B0rk incident from the Scripting for Server Security Part 2 lab. Use the sample report from Google (linked under Resources) as your template.
\end{enumerate}

\section*{Acceptance Criteria}
\begin{itemize}
    \item \textbf{Everything must be cited.} Every command, statement, and line of code in your scripts and documents must be cited, either with a link to a reference or with a note naming the previous class where you learned it, for example: \texttt{\# Learned in CIS 117 Intro to Linux}. In scripts, put the citation in a comment next to the code. Anything without a citation will not count.
    \item The restore is proven with before and after screenshots.
    \item The restored files are checked against the originals, for example with \texttt{sha256sum}.
    \item Your cron schedule is shown working.
    \item Scripts are well commented and uploaded to GitHub.
    \item Every terminal screenshot shows your prompt with your username and hostname, and either the output of \texttt{date} or a note beside the terminal with your name and the term. Hostnames must match your server inventory from the Server Setup lab.
    \item Your GitHub repository shows your work over time: several commits across more than one day, each with a message describing the change, not a single upload of finished code.
\end{itemize}

\section*{Deliverables}
Audience: a teammate who has to restore a server during an emergency. Assume technical knowledge but not expertise.  Screenshots are helpful to go with your descriptions.

Make sure your deliverables include documents for BOTH servers, clearly labeled with what the document is and which server it's for.
\begin{itemize}
    \item Your backup script, with a link to your GitHub where you've uploaded it.
    \item A restore procedure someone else could follow during an emergency, including how long your restore drill took.
    \item Your post mortem.
\end{itemize}

\section*{Resources}
References, a video, a PowerPoint and some notes are available at my website
Link: \href{https://www.aholdengouveia.name/LinuxAdmin/backups.html}{\textbf{\nolinkurl{https://www.aholdengouveia.name/LinuxAdmin/backups.html}}}

My video on how to use tar
Link: \href{https://www.youtube.com/watch?v=x0EK6PmKRCU}{\textbf{\nolinkurl{https://www.youtube.com/watch?v=x0EK6PmKRCU}}}

Sample post mortem report from Google
Link: \href{https://web.archive.org/web/20220724194136/https://developers.googleblog.com/2013/05/google-api-infrastructure-outage_3.html}{\textbf{\nolinkurl{https://web.archive.org/web/20220724194136/https://developers.googleblog.com/2013/05/google-api-infrastructure-outage_3.html}}}

% ============================================================
% EXAMPLE: Accessible Table (if needed)
% Uncomment and modify the lines below to add a table
% This ensures proper screen reader support
% ============================================================
%\begin{table}[h]
%\caption{Command Options}
%\begin{tabular}{|l|l|}
%\hline
%\textbf{Option} & \textbf{Description} \\
%\hline
%-F & Field separator \\
%\hline
%-v & Variable assignment \\
%\hline
%\end{tabular}
%\end{table}

\end{document}